\section{Introducci\'{o}n}
\label{sec:introduction}

% Problema de investigaci\'{o}n
[Describir el problema: ineficiencias en la gesti\'{o}n de inventarios en empresas log\'{i}sticas peruanas. Costos de sobrestock y quiebre de stock. Necesidad de herramientas anal\'{i}ticas avanzadas.]

% Justificaci\'{o}n
[Justificar por qu\'{e} machine learning es una alternativa superior a m\'{e}todos tradicionales (EOQ, modelos determin\'{i}sticos) para este contexto.]

% Objetivo general
El objetivo general de esta investigaci\'{o}n es desarrollar un sistema de anal\'{i}tica prescriptiva basado en machine learning que optimice la gesti\'{o}n de inventarios en una empresa del sector log\'{i}stico peruano.

% Objetivos espec\'{i}ficos
Los objetivos espec\'{i}ficos son:
\begin{enumerate}
    \item Analizar los patrones de demanda y comportamiento de compras mediante an\'{a}lisis exploratorio de datos.
    \item Implementar y comparar m\'{u}ltiples algoritmos de machine learning para la predicci\'{o}n de demanda.
    \item Desarrollar un m\'{o}dulo prescriptivo que genere recomendaciones \'{o}ptimas de reabastecimiento.
    \item Evaluar el desempe\~{n}o del sistema propuesto frente a m\'{e}todos tradicionales de gesti\'{o}n de inventarios.
\end{enumerate}

% Estructura del documento
[Describir la organizaci\'{o}n del resto del documento.]
